% Fixture: presentación beamer con overlays, bloques y columnas.
\documentclass{beamer}
\usepackage[T1]{fontenc}
\usepackage[utf8]{inputenc}
\usepackage[spanish]{babel}
\usetheme{Madrid}

\title{Presentación de prueba}
\author{Fixture de pruebas}
\date{}

\begin{document}

\begin{frame}
  \titlepage
\end{frame}

\begin{frame}{Overlays}
  \begin{itemize}
    \item<1-> Aparece en la primera diapositiva
    \item<2-> Aparece en la segunda
    \item<3-> Aparece en la tercera
  \end{itemize}
\end{frame}

\begin{frame}{Bloques y columnas}
  \begin{columns}
    \begin{column}{0.48\textwidth}
      \begin{block}{Bloque}
        Contenido normal.
      \end{block}
    \end{column}
    \begin{column}{0.48\textwidth}
      \begin{alertblock}{Alerta}
        $a^2 + b^2 = c^2$
      \end{alertblock}
    \end{column}
  \end{columns}
\end{frame}

\end{document}
